
\documentclass[a4paper]{article} %
\usepackage{graphicx,amssymb} %

\usepackage{color}

\textwidth=15cm \hoffset=-1.2cm %
\textheight=25cm \voffset=-2cm %

\pagestyle{empty} %

\date{} %

\def\keywords#1{\begin{center}{\bf Keywords}\\{#1}\end{center}} %

\def\titulo#1{\title{#1}} %
\def\autores#1{\author{#1}} %

% Please, do not change any of the above lines


\begin{document}

% Type down your paper title
\titulo{Population--based anamorphosis maps\\ for railway radial networks}

% Authors
\autores{Eugenio Roanes--Lozano \\ %
         Instituto de Matem\'atica Interdisciplinar (IMI),\\
         Algebra Dept., Universidad Complutense de Madrid, Spain\\ \\ % Affiliation 1
         Alberto Garc\'{\i}a--\'Alvarez \\ %
         Deputy Director for Renfe (Spanish Railways) Passengers Services, Spain \\ \\ % Affiliation 2
         Jos\'e Luis Gal\'an--Garc\'{\i}a \\ %
         Applied Mathematics Dept., Universidad de M\'alaga, Spain \\ \\ % Affiliation 3
         Luis Mesa \\ %
         Spanish Railways Foundation, Spain\\ \\ % Affiliation 4
         \tt{eroanes@mat.ucm.es} % Only one corresponding e-mail
       }%


\maketitle

\thispagestyle{empty}



% The abstract

\begin{abstract}
The Spanish railway network is radial but very complex to
operate, because two different track gauges, five signalling
systems and two electrification systems coexist. Therefore, how
to go on developing the high speed network and which are the best
routes for trains are complicated questions. We are developing,
in cooperation with the Spanish Railway Foundation, software
packages that can be aids to decision making in these two issues.
Yet one more step in this direction is presented here.
\end{abstract}

\keywords{Radial railway networks, Anamorphosis maps, Isochrone
circle maps, Computer algebra systems}
% Write down at least 3 Keywords

\section{Introduction}

The Spanish railway network is radial but very complex to
operate, because:
\begin{itemize}
\item two different track gauges coexist: the so
      called \textit{Iberian} gauge ($1667 mm$) and the
      \textit{international gauge} ($1435 mm$) (used in the high speed
      network),
\item two electrification systems have been used (3000 V DC,
      25000 V AC) and, finally,
\item there are different signalling systems
      (ASFA, ASFA 200, LZB, EBICAB, ERTMS).
\end{itemize}
The shape of the network is not due to a katabasis or anabasis of
the whole Spanish society, but to the location of Madrid (the
biggest city and capital) in the centre of the country and the
location of the rest of big cities in the periphery. In 2013 the
fares have been lowered and different discounts (for instance in
low demand periods) are offered in order to attract more
travelers and to increase the fraction of population using the
high-speed trains.

Regarding the two gauges, there are gauge changeovers at several
points, so that both subnetworks are connected. A subset of the
rolling stock is dual gauge. Regarding electrification and
signalling systems, many locomotives and multiple units can read
different signalling systems, are multi-voltage. Even hybrid
rolling stock has been developed (730 series).

The high speed network has grown very quickly, and only China has
nowadays a longer high speed railway network. All new lines have
been built with double track and top technologies ($\geq 300
km/h$ track design, \textit{ERTMS} traffic management system,
$25000 KV$ AC electrification, etc.). The growth of the network
has been supported by the different governments and has only been
slowed down due to the economical crisis.

Nevertheless, how to go on developing the high speed network and
which are the best routes for trains are complicated questions.
We are developing, in cooperation with the \textit{Spanish
Railway Foundation}, software packages that can be aids to
decision making in these two issues. We have followed two lines:
\begin{itemize}
\item We have developed (within the frame of two research projects signed between the \textit{Spanish
      Railways Foundation} and the \textit{Universidad Complutense de Madrid} and the \textit{Universidad
      Polit\'ecnica de Madrid}) a computer package that is able to calculate precise timings, consumptions,
      costs, emissions, best routes, etc., for each piece of \textit{Renfe}'s (main railway operator)
      rolling stock running on \textit{Adif}'s (infrastructure company) lines [1].
\item We have also developed what we have called \textit{isochrone circle graphs} and a
      \textit{geometric index} for radial railway networks improvement estimation [2].
      \textit{Isochrone circle graphs} were inspired
      by pie charts (also known as circle graphs), polar area diagrams (similar to usual pie charts, but
      sectors are equal angles and their area is adjusted changing their radii instead of their amplitude)
      and anamorphosis maps (also known as central point cartograms or distance cartograms;
      where the geometry of the country or region is distorted according to the time that it takes
      to travel to different peripheral destinations from a central origin). An \textit{isochrone circle
      graph} corresponding to Spain in 2013 (centred at Madrid) can be found in the figure below.

%***************** FIGURA **************
\begin{figure}[htb]
\begin{center}
\scalebox{.65}{
\includegraphics%[width=\textwidth]
{Fig8.jpg}
} %
%\caption{\textit{Isochrone circle graph}, reflecting 2013\newline
%         population and timings in the Spanish railway network.}
\end{center}
\end{figure}
%***************************************

      We have followed two approaches to compute and draw \textit{isochrone circle graphs}:
      \begin{itemize}
      \item The first approach [2] was illustrated with a sketch constructed with a Dynamic Geometry System and
            used sliders to change the input parameters (timing to each peripheral destination and population
            of these destinations). It was very comfortable to use, but the number of destinations
            considered was somehow fixed (changing it required to construct a complete new sketch).
      \item In the second approach we designed and implemented a complete new package in the CAS \textit{Maple}
            that takes as input the lists of destinations, best timings and populations and builds the
            corresponding \textit{isochrone circle graphs} and performs all the corresponding calculations [3].
            This approach has yet another advantage: symbolic computations can be performed, and therefore
            parameters can be introduced in the computations.
       \end{itemize}
\end{itemize}

An improved version of [3] will be presented here. It has to be
emphasized that now a population--based version of an
anamorphosis map can be drawn or superimposed to the
\textit{isochrone circle graph}.


\bibliographystyle{plain}
\begin{thebibliography}{10}
\bibitem{CiSE}
A. Hernando, E. Roanes--Lozano, A. Garc\'{\i}a--\'Alvarez, L. Mesa, I. Gonz\'alez--Franco:
Optimal Route Finding and Rolling--Stock Selection for the Spanish Railways.
\textit{Comp. in Sci. \& Eng.} 14/4 (2012) 82--89.
DOI: http://dx.doi.org/10.1109/MCSE.2012.80.

\bibitem{RACSAM}
E. Roanes--Lozano, A. Garc\'{\i}a--\'Alvarez, A. Hernando:
A geometric approach to the estimation of radial railway network improvement.
\textit{Rev. R. Acad. Cienc. Exactas F\'is. Nat., Ser. A Mat. RACSAM}
106/1 (2012) 35--46.
DOI: http://dx.doi.org/10.1007/s13398-011-0050-6.

\bibitem{FEMTEC}
E. Roanes--Lozano, A. Garc\'{\i}a--\'Alvarez, J. L. Gal\'an--Garc\'{\i}a, L. Mesa:
Estimating radial railway network improvement with a CAS.
To be presented at:
\textit{FEMTEC'2013 (4th International Congress on Computational Engineering and
         Sciences)}, Las Vegas, May 19--24, 2013.

\end{thebibliography}


\end{document}
